\documentclass[10pt,a4paper]{amsart}
\usepackage{amsmath,amssymb,amsthm,graphicx}
% For commutative diagrams we use Xy-pic
\usepackage[all]{xy}

% We use 2cell for 2-commutative diagrams.
\xyoption{2cell}
% Theorem environments.
%
\theoremstyle{plain}
\newtheorem{theorem}[subsection]{Theorem}
\newtheorem{proposition}[subsection]{Proposition}
\newtheorem{lemma}[subsection]{Lemma}

\theoremstyle{definition}
\newtheorem{definition}[subsection]{Definition}
\newtheorem{example}[subsection]{Example}
\newtheorem{exercise}[subsection]{Exercise}
\newtheorem{situation}[subsection]{Situation}

\theoremstyle{remark}
\newtheorem{remark}[subsection]{Remark}
\newtheorem{remarks}[subsection]{Remarks}

\numberwithin{equation}{subsection}

% Macros
%
\def\lim{\mathop{\mathrm{lim}}\nolimits}
\def\colim{\mathop{\mathrm{colim}}\nolimits}
\def\Spec{\mathop{\mathrm{Spec}}}
\def\Hom{\mathop{\mathrm{Hom}}\nolimits}
\def\Ext{\mathop{\mathrm{Ext}}\nolimits}
\def\SheafHom{\mathop{\mathcal{H}\!\mathit{om}}\nolimits}
\def\SheafExt{\mathop{\mathcal{E}\!\mathit{xt}}\nolimits}
\def\Sch{\mathit{Sch}}
\def\Mor{\mathop{\mathrm{Mor}}\nolimits}
\def\Ob{\mathop{\mathrm{Ob}}\nolimits}
\def\Sh{\mathop{\mathit{Sh}}\nolimits}
\def\NL{\mathop{N\!L}\nolimits}
\def\CH{\mathop{\mathrm{CH}}\nolimits}
\def\proetale{{pro\text{-}\acute{e}tale}}
\def\etale{{\acute{e}tale}}
\def\QCoh{\mathit{QCoh}}
\def\Ker{\mathop{\mathrm{Ker}}}
\def\Im{\mathop{\mathrm{Im}}}
\def\Coker{\mathop{\mathrm{Coker}}}
\def\Coim{\mathop{\mathrm{Coim}}}

% Boxtimes
%
\DeclareMathSymbol{\boxtimes}{\mathbin}{AMSa}{"02}

%
% Macros for moduli stacks/spaces
%
\def\QCohstack{\mathcal{QC}\!\mathit{oh}}
\def\Cohstack{\mathcal{C}\!\mathit{oh}}
\def\Spacesstack{\mathcal{S}\!\mathit{paces}}
\def\Quotfunctor{\mathrm{Quot}}
\def\Hilbfunctor{\mathrm{Hilb}}
\def\Curvesstack{\mathcal{C}\!\mathit{urves}}
\def\Polarizedstack{\mathcal{P}\!\mathit{olarized}}
\def\Complexesstack{\mathcal{C}\!\mathit{omplexes}}
% \Pic is the operator that assigns to X its picard group, usage \Pic(X)
% \Picardstack_{X/B} denotes the Picard stack of X over B
% \Picardfunctor_{X/B} denotes the Picard functor of X over B
\def\Pic{\mathop{\mathrm{Pic}}\nolimits}
\def\Picardstack{\mathcal{P}\!\mathit{ic}}
\def\Picardfunctor{\mathrm{Pic}}
\def\Deformationcategory{\mathcal{D}\!\mathit{ef}}
\usepackage{cleveref}
\crefname{originallemma}{original lemma}{original lemmas}
\Crefname{originallemma}{Original lemma}{Original lemmas}
\makeatletter
\newlabel{algebra-lemma-support-point}{{00L5}{source}{}{sourcepointer-1}{}}
\newlabel{algebra-lemma-support-point@cref}{{[originallemma][0][]00L5}{[1][0][]source}{}{}{}}
\makeatother
\title{Existence of a generating sequence whose successive quotients are the residue field at a minimal prime of the support of a finite module}
\author{The Stacks Project Authors}
\date{stacks-project:chow.tex:local-snapshot-20261012; GNU FDL 1.2 or later}
\begin{document}
\maketitle
\noindent\textit{Editorial note.} All delivered mathematical text is the author's own, copied from the pinned Stacks Project chapter file chow.tex; no statement or proof step was rewritten, reordered or removed.\par\medskip
\noindent\textit{Editorial note.} Preamble provenance: the pinned chapter file opens with an input of the shared preamble and carries no preamble text of its own, so the project's single Stacks-wide preamble staged at sources/stacks-chow/preamble.tex is excerpted verbatim as two spans, lines 14-18 (Xy-pic setup) and lines 154-217 (the native theorem declarations and mathematical macros). Nothing was added, removed or redefined; the wrapper class and the amsmath/amssymb/amsthm/graphicx base packages are supplied by the assembler.\par\medskip
\noindent\textit{Editorial note.} The counted statement and proof contain no cross reference, no citation and no included file, so the delivered body is self contained.\par\medskip
\noindent\textit{Editorial note.} Canonical location of the counted proposition: Stacks tag 02QC (sources/stacks-chow/pinned-original-source.tex line 17964).\par\medskip
\noindent\textit{Allowed standard prerequisites (editorial).}\begin{itemize}
\item Support of a finite module over a Noetherian local ring and finiteness of the length at a minimal prime of the support. Algebra, Lemma labelled algebra-lemma-support-point (Stacks tag 00L5), invoked at proof line 17982. Supplies that q is a minimal prime of the support exactly when M\_q has finite length, which is what makes l = length\_\{R\_q\}(M\_q) finite and the sequence y\_1, ..., y\_l available; the lifting and clearing of denominators is inside the counted span.
\end{itemize}\medskip
\ifdefined\hypertarget\hypertarget{sourcepointer-1}{}\fi
\noindent\textit{Source reference.} Original lemma 00L5: sources/stacks-chow/tags.txt, source lines 775--775. Cross-chapter standard prerequisite located in the project tag list staged in the local snapshot; invoked at proof line 17982 of the counted proof. It is retained only as a navigation pointer; every load-bearing step is inside the delivered span.\par\medskip
\begin{lemma}
\label{lemma-good-sequence-exists}
Let $R$ be a local Noetherian ring.
Let $\mathfrak q \subset R$ be a prime ideal.
Let $M$ be a finite $R$-module such that
$\mathfrak q$ is one of the minimal primes of the support of $M$.
Then there exist $x_1, \ldots, x_l \in M$ such that
\begin{enumerate}
\item the support of $M / \langle x_1, \ldots, x_l\rangle$ does not contain
$\mathfrak q$, and
\item $\langle x_1, \ldots, x_i\rangle / \langle x_1, \ldots, x_{i - 1}\rangle
\cong R/\mathfrak q$ for $i = 1, \ldots, l$.
\end{enumerate}
Moreover, in this case $l = \text{length}_{R_\mathfrak q}(M_\mathfrak q)$.
\end{lemma}

\begin{proof}
The condition that $\mathfrak q$ is a minimal prime in the support
of $M$ implies that $l = \text{length}_{R_\mathfrak q}(M_\mathfrak q)$
is finite (see Algebra, Lemma \ref{algebra-lemma-support-point}).
Hence we can find $y_1, \ldots, y_l \in M_{\mathfrak q}$
such that
$\langle y_1, \ldots, y_i\rangle / \langle y_1, \ldots, y_{i - 1}\rangle
\cong \kappa(\mathfrak q)$ for $i = 1, \ldots, l$.
We can find $f_i \in R$, $f_i \not \in \mathfrak q$ such that
$f_i y_i$ is the image of some element $z_i \in M$.
Moreover, as $R$ is Noetherian we can write
$\mathfrak q = (g_1, \ldots, g_t)$ for some $g_j \in R$.
By assumption $g_j y_i \in \langle y_1, \ldots, y_{i - 1} \rangle$
inside the module $M_{\mathfrak q}$.
By our choice of $z_i$ we can find some further elements
$f_{ji} \in R$, $f_{ij} \not \in \mathfrak q$ such that
$f_{ij} g_j z_i \in \langle z_1, \ldots, z_{i - 1} \rangle$
(equality in the module $M$).
The lemma follows by taking
$$
x_1 = f_{11}f_{12}\ldots f_{1t}z_1,
\quad
x_2 = f_{11}f_{12}\ldots f_{1t}f_{21}f_{22}\ldots f_{2t}z_2,
$$
and so on. Namely, since all the elements $f_i, f_{ij}$ are invertible
in $R_{\mathfrak q}$ we still have that
$R_{\mathfrak q}x_1 + \ldots + R_{\mathfrak q}x_i /
R_{\mathfrak q}x_1 + \ldots + R_{\mathfrak q}x_{i - 1}
\cong \kappa(\mathfrak q)$ for $i = 1, \ldots, l$.
By construction, $\mathfrak q x_i \in \langle x_1, \ldots, x_{i - 1}\rangle$.
Thus $\langle x_1, \ldots, x_i\rangle / \langle x_1, \ldots, x_{i - 1}\rangle$
is an $R$-module generated by one element, annihilated $\mathfrak q$
such that localizing at $\mathfrak q$ gives a $q$-dimensional
vector space over $\kappa(\mathfrak q)$.
Hence it is isomorphic to $R/\mathfrak q$.
\end{proof}

\end{document}
